% =====================================================================
%  整理者　macanine <macanine@163.com>
%  仓库　　github.com/macanine/zju-linear-algebra-papers
%  来源　　解析据 papers/2021-2022-秋冬-期中-甲.tex 撰写（原卷见 sources/2016-2025-期中-甲-历年真题汇编.pdf 第 10 页）
% =====================================================================

\begin{solution}
记所给 $n$ 阶行列式为 $D_n$（并约定 $D_0=1$，即空行列式的值）。

\textbf{(1)} 按定义直接算：
\[
D_1=|7|=7,\qquad
D_2=\begin{vmatrix}7&-6\\-1&7\end{vmatrix}=7\times7-(-6)\times(-1)=43 .
\]

\textbf{(2)} 先找递推式。按第一列展开，第一列只有 $a_{11}=7$ 与 $a_{21}=-1$ 非零：
\[
D_n=7M_{11}+M_{21}=7D_{n-1}+M_{21},
\]
其中 $M_{11}=D_{n-1}$；而删去第 $2$ 行、第 $1$ 列后，余下矩阵的第一行是 $(-6,0,\dots,0)$，
按这一行展开得 $M_{21}=(-6)\cdot D_{n-2}$（余下的正是同形状的 $n-2$ 阶行列式）。于是
\[
D_n=7D_{n-1}-6D_{n-2}\qquad(n\ge2).
\]

再作归纳。$n=0,1$ 时公式 $(6^{n+1}-1)/5$ 成立（分别为 $1$ 与 $7$）。
设 $n\le k$ 时 $D_n=\dfrac{6^{n+1}-1}{5}$，则特别地 $D_k=\dfrac{6^{k+1}-1}{5}$，$D_{k-1}=\dfrac{6^{k}-1}{5}$，代入递推式
\[
\begin{aligned}
D_{k+1}&=7D_k-6D_{k-1}
=7\cdot\frac{6^{k+1}-1}{5}-6\cdot\frac{6^{k}-1}{5}\\
&=\frac{7\cdot6^{k+1}-7-6^{k+1}+6}{5}
=\frac{6\cdot6^{k+1}-1}{5}
=\frac{6^{k+2}-1}{5}.
\end{aligned}
\]
代 $n=3$ 检验：$D_3=7\times43-6\times7=259=\dfrac{6^{4}-1}{5}$。
\ans{$D_1=7$，$D_2=43$；且 $D_n=\dfrac{6^{n+1}-1}{5}$，故 $D_{k+1}=\dfrac{6^{k+2}-1}{5}$。}
\end{solution}

\begin{solution}
由 $|B|=1\ne0$ 知 $B$ 可逆，于是 $BA=B^{T}B$ 给出
\[
A=B^{-1}B^{T}B .
\]
取行列式，注意 $|B^{T}|=|B|$、$|B^{-1}|=|B|^{-1}$：
\[
|A|=|B^{-1}|\cdot|B^{T}|\cdot|B|=|B|=1 .
\]
故 $A$ 可逆，且 $A^{*}=|A|A^{-1}=A^{-1}$。于是
\[
|A^{*}-2E|=|A^{-1}-2E|=|A^{-1}|\cdot|E-2A|=|E-2A| .
\]
把 $A=B^{-1}B^{T}B$ 代入，注意 $E=B^{-1}B$：
\[
E-2A=B^{-1}B-2B^{-1}B^{T}B=B^{-1}\bigl(B-2B^{T}B\bigr),
\]
两边取行列式并用 $|B|=1$，得 $|E-2A|=|B-2B^{T}B|$。逐项算出
\[
B^{T}B=\begin{pmatrix}1&1&1\\1&2&2\\1&2&3\end{pmatrix},\qquad
B-2B^{T}B=\begin{pmatrix}-1&-1&-1\\-2&-3&-3\\-2&-4&-5\end{pmatrix}.
\]
对右边矩阵作行倍加（不改变行列式）：$R_2-2R_1$，$R_3-2R_1$，再 $R_3-2R_2$，化为上三角
\[
\begin{pmatrix}-1&-1&-1\\0&-1&-1\\0&0&-1\end{pmatrix},
\]
故 $|B-2B^{T}B|=(-1)^{3}=-1$，从而 $|A^{*}-2E|=-1$。

顺带算出
\[
A=B^{-1}B^{T}B=\begin{pmatrix}0&-1&-1\\0&0&-1\\1&2&3\end{pmatrix},\qquad |A|=1,
\]
与上面一致。
\ans{$|A^{*}-2E|=-1$。}
\rem{这里 $A$ 是三阶方阵，所以 $A^{*}=|A|A^{-1}$、进而 $A^{*}=A^{-1}$ 成立的前提是 $|A|\ne0$；
本题由 $|A|=|B|=1$ 保证。若误把 $A^{*}$ 与 $A$ 的阶数关系写成 $|A^{*}|=|A|^{n}$，结果就会差一个因子。}
\end{solution}

\begin{solution}
由 $A^{2}=A$ 得
\[
A(A-E)=A^{2}-A=O .
\]
用到两条关于秩的基本事实，设 $X,Y$ 为 $n$ 阶方阵：

(i) 若 $XY=O$，则 $\operatorname{Im}Y\subseteq\ker X$，由秩－零化度定理
$r(Y)\le\dim\ker X=n-r(X)$，即 $r(X)+r(Y)\le n$；

(ii) $\operatorname{Im}(X+Y)\subseteq\operatorname{Im}X+\operatorname{Im}Y$，故 $r(X+Y)\le r(X)+r(Y)$。

\textbf{上界。}在 (i) 中取 $X=A$，$Y=A-E$，由 $A(A-E)=O$ 得
\[
r(A)+r(A-E)\le n .
\]

\textbf{下界。}在 (ii) 中取 $X=A$，$Y=E-A$，注意 $A+(E-A)=E$ 且 $r(E-A)=r(A-E)$：
\[
n=r(E)=r\bigl(A+(E-A)\bigr)\le r(A)+r(E-A)=r(A)+r(A-E).
\]

两式合并即得 $r(A)+r(A-E)=n$。
\qed
\end{solution}

\begin{solution}
\textbf{(1)} 按未知量 $x_1,x_2,x_3,x_4$ 的顺序取出系数：
\[
A=\begin{pmatrix}1&1&-1&2\\2&3&-3&5\\3&1&-2&3\\2&1&-5&-1\end{pmatrix}.
\]

\textbf{(2)} 对 $A$ 作初等行变换：先以第一行消去第一列其余元素，
\[
A\to\begin{pmatrix}1&1&-1&2\\0&1&-1&1\\0&-2&1&-3\\0&-1&-3&-5\end{pmatrix}
\qquad\begin{aligned}
R_2&\leftarrow R_2-2R_1,\\
R_3&\leftarrow R_3-3R_1,\\
R_4&\leftarrow R_4-2R_1,
\end{aligned}
\]
再以第二行消去第二列下方元素，
\[
\to\begin{pmatrix}1&1&-1&2\\0&1&-1&1\\0&0&-1&-1\\0&0&-4&-4\end{pmatrix}
\qquad\begin{aligned}
R_3&\leftarrow R_3+2R_2,\\
R_4&\leftarrow R_4+R_2,
\end{aligned}
\]
最后 $R_4\leftarrow R_4-4R_3$，得阶梯形
\[
U=\begin{pmatrix}1&1&-1&2\\0&1&-1&1\\0&0&-1&-1\\0&0&0&0\end{pmatrix}.
\]

\textbf{(3)} $U$ 有三个非零行（主元分别在 $1,2,3$ 列），故 $r(A)=r(U)=3$。

\textbf{(4)} 继续化为行最简形：$R_3\leftarrow-R_3$ 后依次作 $R_2+R_3$，$R_1+R_3$，$R_1-R_2$，得
\[
\begin{pmatrix}1&0&0&1\\0&1&0&2\\0&0&1&1\\0&0&0&0\end{pmatrix}.
\]
同解方程组为 $x_1=-x_4$，$x_2=-2x_4$，$x_3=-x_4$。未知量 $4$ 个、秩为 $3$，故自由未知量恰一个 $x_4$；
令 $x_4=t$（$t$ 为任意实数）即得通解。代回原方程组四个方程均成立。
\ans{$A=\begin{pmatrix}1&1&-1&2\\2&3&-3&5\\3&1&-2&3\\2&1&-5&-1\end{pmatrix}$，
$U=\begin{pmatrix}1&1&-1&2\\0&1&-1&1\\0&0&-1&-1\\0&0&0&0\end{pmatrix}$，$r(A)=3$，
通解为 $X=t\begin{pmatrix}-1\\-2\\-1\\1\end{pmatrix}$（$t\in\mathbb{R}$）。}
\end{solution}

\begin{solution}
先定 $|A|$。$A^{*}$ 是主对角线全为 $1$ 的下三角阵，故 $|A^{*}|=1$。
又对 $n$ 阶方阵有 $|A^{*}|=|A|^{n-1}$，这里 $n=3$，所以
\[
|A|^{2}=1\ \Longrightarrow\ |A|=\pm1 .
\]
题目给出 $|A|>0$，故 $|A|=1$，从而 $A$ 可逆且
\[
A^{*}=|A|A^{-1}=A^{-1},\qquad\text{即}\qquad A^{-1}=A^{*}.
\]
由 $AB=E+3A$，两边左乘 $A^{-1}$：
\[
B=A^{-1}(E+3A)=A^{-1}+3E=A^{*}+3E
=\begin{pmatrix}1&0&0\\0&1&0\\1&0&1\end{pmatrix}+3E
=\begin{pmatrix}4&0&0\\0&4&0\\1&0&4\end{pmatrix}.
\]
回代验算：$AB=A(A^{*}+3E)=AA^{*}+3A=|A|E+3A=E+3A$，满足方程。顺带算出
\[
A=(A^{*})^{-1}=\begin{pmatrix}1&0&0\\0&1&0\\-1&0&1\end{pmatrix},\qquad |A|=1>0,
\]
与题设 $|A|>0$ 相容。
\ans{$B=\begin{pmatrix}4&0&0\\0&4&0\\1&0&4\end{pmatrix}$。}
\end{solution}

\begin{solution}
记
\[
N=A-3E=\begin{pmatrix}0&1&0\\0&0&1\\0&0&0\end{pmatrix},
\]
则 $A=3E+N$，其中 $3E$ 与 $N$ 可交换，且
\[
N^{2}=\begin{pmatrix}0&0&1\\0&0&0\\0&0&0\end{pmatrix}\ne O,\qquad N^{3}=O .
\]
$N$ 幂零，故二项式定理可用：对 $k\ge3$
\[
A^{k}=(3E+N)^{k}=\sum_{i=0}^{k}\binom{k}{i}3^{k-i}N^{i}
=3^{k}E+k\,3^{k-1}N+\binom{k}{2}3^{k-2}N^{2},
\]
求和只到 $i=2$ 是因为 $N^{3}=O$。由此逐项求值。

\textbf{(1)} 取 $k=2$：
\[
A^{2}=9E+6N+N^{2}=9E+6\begin{pmatrix}0&1&0\\0&0&1\\0&0&0\end{pmatrix}+\begin{pmatrix}0&0&1\\0&0&0\\0&0&0\end{pmatrix}
=\begin{pmatrix}9&6&1\\0&9&6\\0&0&9\end{pmatrix}.
\]

\textbf{(2)} 取 $k=3$：
\[
A^{3}=27E+27N+9N^{2}=\begin{pmatrix}27&27&9\\0&27&27\\0&0&27\end{pmatrix}.
\]

\textbf{(3)} 取 $k=100$，$\dbinom{100}{2}=\dfrac{100\times99}{2}=4950$：
\[
A^{100}=3^{100}E+100\cdot3^{99}N+4950\cdot3^{98}N^{2}
=\begin{pmatrix}3^{100}&100\cdot3^{99}&4950\cdot3^{98}\\0&3^{100}&100\cdot3^{99}\\0&0&3^{100}\end{pmatrix}.
\]
（与 (2) 对照：把 $k=3$ 代入 $A^{k}$ 的公式即得 $A^{3}$，$27,27,9$ 分别对应 $3^{3},3\cdot3^{2},\binom{3}{2}3^{1}$。）

\textbf{(4)} 由 $A=3\Bigl(E+\dfrac{N}{3}\Bigr)$ 及 $N^{3}=O$，对 $X=\dfrac{N}{3}$ 用 $(E+X)^{-1}=E-X+X^{2}$（因 $E=(E+X)(E-X+X^{2})$）：
\[
A^{-1}=\frac13\Bigl(E-\frac{N}{3}+\frac{N^{2}}{9}\Bigr)
=\frac13E-\frac19N+\frac1{27}N^{2}
=\frac{1}{27}\begin{pmatrix}9&-3&1\\0&9&-3\\0&0&9\end{pmatrix}.
\]
把 $A$ 乘回即可验证 $A^{-1}A=E$：$N,N^{2},N^{3}$ 的系数两两抵消。

\textbf{(5)} 取
\[
f(x)=(x-3)^{3}=x^{3}-9x^{2}+27x-27 ,
\]
即 $a=-9$，$b=27$，$c=-27$，则
\[
f(A)=(A-3E)^{3}=N^{3}=O .
\]
再说明次数不能再低：$A$ 的特征多项式为 $(x-3)^{3}$，极小多项式与特征多项式有相同的根，
故极小多项式必为 $(x-3)^{k}$（$k=1,2,3$）；而 $(A-3E)^{2}=N^{2}\ne O$、$A-3E=N\ne O$，
所以 $k=3$，$(x-3)^{3}$ 正是极小多项式，它是零化 $A$ 的最低次多项式。
\ans{(1) $A^{2}=\begin{pmatrix}9&6&1\\0&9&6\\0&0&9\end{pmatrix}$；
(2) $A^{3}=\begin{pmatrix}27&27&9\\0&27&27\\0&0&27\end{pmatrix}$；
(3) $A^{100}=\begin{pmatrix}3^{100}&100\cdot3^{99}&4950\cdot3^{98}\\0&3^{100}&100\cdot3^{99}\\0&0&3^{100}\end{pmatrix}$；
(4) $A^{-1}=\dfrac{1}{27}\begin{pmatrix}9&-3&1\\0&9&-3\\0&0&9\end{pmatrix}$；
(5) $f(x)=(x-3)^{3}=x^{3}-9x^{2}+27x-27$（即 $a=-9$，$b=27$，$c=-27$）。}
\end{solution}

\begin{solution}
把元素全部模 $2$ 来看。记 $\bar{a}$ 为整数 $a$ 模 $2$ 的余数，$\bar A=(\bar a_{ij})$。当 $i>j$ 时
\[
i+j\ge2+1=3\ \Longrightarrow\ (i+j)!\ \text{含因子}\ 2,
\]
故 $\bar a_{ij}=0$；当 $i=j$ 时 $a_{ii}=2021=2\times1010+1$ 为奇数，故 $\bar a_{ii}=1$。于是 $\bar A$ 是
主对角线全为 $1$ 的上三角矩阵，从而 $\det\bar A=1$。

行列式是各元素的整系数多项式，取模 $2$ 与取行列式可交换，即
\[
\overline{|A|}\equiv\det\bar A=1\pmod 2 .
\]
所以 $|A|$ 是奇数，特别地 $|A|\ne0$。

（可任取 $n$ 检验：$n=2$ 时 $|A|=4084435$，为奇数。）
\qed
\end{solution}

\begin{solution}
由矩阵的标准形理论，$\operatorname{rank}A=r$ 意味存在可逆矩阵 $P,Q$ 使
\[
PAQ=J:=\begin{pmatrix}E_r&0\\0&0\end{pmatrix}.
\]
（把 $A$ 经初等行、列变换化为标准形即可得到这样的 $P,Q$，$J$ 的左上角是 $r$ 阶单位阵。）
由此 $A=P^{-1}JQ^{-1}$。

\textbf{构造。}取
\[
B=QJQ^{-1},\qquad C=P^{-1}JP .
\]
$B$ 是 $J$ 的相似变换、$C$ 也是（$C=P^{-1}JP$），乘可逆矩阵不改变秩，而 $r(J)=r$，故
\[
r(B)=r(J)=r,\qquad r(C)=r(J)=r,
\]
且 $B,C$ 都是 $n$ 阶实方阵，符合要求。

\textbf{验证 $AB=CA$。}注意 $J^{2}=J$（$J$ 是幂等阵），且 $PP^{-1}=QQ^{-1}=E$：
\[
AB=P^{-1}JQ^{-1}\cdot QJQ^{-1}=P^{-1}J^{2}Q^{-1}=P^{-1}JQ^{-1}=A,
\]
\[
CA=P^{-1}JP\cdot P^{-1}JQ^{-1}=P^{-1}J^{2}Q^{-1}=P^{-1}JQ^{-1}=A .
\]
两个乘积都等于 $A$ 本身，当然相等，即 $AB=CA$。故所求的 $B,C$ 存在。
\qed
\rem{构造的要害是让 $AB$ 与 $CA$ 都等于 $A$：把 $A$ 写成 $P^{-1}JQ^{-1}$ 后，分别用 $QJQ^{-1}$、$P^{-1}JP$ 去“抵消”两侧的因子。}
\end{solution}
